%-------------------------
% Resume in Latex
% Author : Jake Gutierrez (template), content: Vrunda Joshi
% Based off of: https://github.com/sb2nov/resume
% License : MIT
%------------------------

\documentclass[letterpaper,11pt]{article}

\usepackage{latexsym}
\usepackage[empty]{fullpage}
\usepackage{titlesec}
\usepackage{marvosym}
\usepackage[usenames,dvipsnames]{color}
\usepackage{verbatim}
\usepackage{enumitem}
\usepackage[hidelinks]{hyperref}
\usepackage{fancyhdr}
\usepackage[english]{babel}
\usepackage{tabularx}
\usepackage{fontawesome5}
\input{glyphtounicode}

%----------FONT OPTIONS----------
% sans-serif
% \usepackage[sfdefault]{FiraSans}
% \usepackage[sfdefault]{roboto}
% \usepackage[sfdefault]{noto-sans}
% \usepackage[default]{sourcesanspro}

% serif
% \usepackage{CormorantGaramond}
% \usepackage{charter}

\pagestyle{fancy}
\fancyhf{} % clear all header and footer fields
\fancyfoot{}
\renewcommand{\headrulewidth}{0pt}
\renewcommand{\footrulewidth}{0pt}

% Adjust margins
\usepackage{geometry}
\geometry{letterpaper, left=1.2cm, top=0.5cm, right=1.2cm, bottom=0.5cm}

\urlstyle{same}

\raggedbottom
\setlength{\parindent}{0pt}
\setlength{\tabcolsep}{0in}

% Sections formatting
\titleformat{\section}{
  \vspace{-4pt}\scshape\raggedright\large
}{}{0em}{}[\color{black}\titlerule \vspace{-5pt}]

% Ensure that generate pdf is machine readable/ATS parsable
\pdfgentounicode=1

%-------------------------
% Custom commands
\newcommand{\resumeItem}[1]{
  \item\small{
    {#1 \vspace{-2pt}}
  }
}

\newcommand{\resumeSubheading}[4]{
  \vspace{-2pt}\item
    \begin{tabularx}{\textwidth}[t]{@{}Xr@{}}
      \textbf{#1} & #2 \\
      \textit{\small#3} & \textit{\small #4} \\
    \end{tabularx}\vspace{-7pt}
}

\newcommand{\resumeSubSubheading}[2]{
    \item
    \begin{tabularx}{\textwidth}{@{}Xr@{}}
      \textit{\small#1} & \textit{\small #2} \\
    \end{tabularx}\vspace{-7pt}
}

\newcommand{\resumeProjectHeading}[2]{
    \vspace{-2pt}\item
    \begin{tabularx}{\textwidth}{@{}Xr@{}}
      \small#1 & #2 \\
    \end{tabularx}\vspace{-7pt}
}

\newcommand{\resumeSubItem}[1]{\resumeItem{#1}\vspace{-4pt}}

\renewcommand\labelitemii{$\vcenter{\hbox{\tiny$\bullet$}}$}

\newcommand{\resumeSubHeadingListStart}{\begin{itemize}[leftmargin=0in, label={}]}
\newcommand{\resumeSubHeadingListEnd}{\end{itemize}}
\newcommand{\resumeItemListStart}{\begin{itemize}[leftmargin=0.15in, label=\textbullet]}
\newcommand{\resumeItemListEnd}{\end{itemize}\vspace{-2pt}}

%-------------------------------------------
%%%%%%  RESUME STARTS HERE  %%%%%%%%%%%%%%%%%%%%%%%%%%%%

\begin{document}

%----------HEADING----------
\begin{center}
    \textbf{\Huge \scshape Vrunda Joshi} \\ \vspace{1pt}
    \small Surat, Gujarat $|$
    \href{mailto:vrundajoshi205@gmail.com}{\underline{vrundajoshi205@gmail.com}} $|$
    \href{https://www.linkedin.com/in/vrunda-joshi-a982b037b}{\underline{linkedin.com/in/vrunda-joshi}} $|$
    \href{https://github.com/VrundaJoshi25}{\underline{github.com/VrundaJoshi25}}
\end{center}

%-----------SUMMARY-----------
\section{Summary}
\small{Mathematics graduate pursuing an M.Sc. in Data Science at Dhirubhai Ambani University, with a strong foundation in statistics, probability, and analytical problem-solving. Skilled in Python, SQL, and machine learning through academic projects and research at SARAL Lab; seeking an entry-level Data Analyst / Data Science role.}

%-----------EXPERIENCE-----------
\section{Experience}
  \resumeSubHeadingListStart
    \resumeSubheading
      {SARAL Lab (Systems, Autonomy, Robotics \& Automation Lab), DAU}{May 2026 - Present}
      {Student Researcher}{Gandhinagar, Gujarat}
      \resumeItemListStart
        \resumeItem{Extracted trip logs from the lab's physics-based simulator of a working coal mine, where self-driving haul trucks navigate a road network of about 240 intersections with every second of every trip recorded.}
        \resumeItem{Visualized raw traffic patterns and preprocessed the trip data: cleaning, structuring, and preparing features to make it modeling-ready.}
        \resumeItem{Trained and benchmarked trip-duration models; the best (Random Forest) was off by just 16 seconds on average, about 7\% of a typical trip, against baselines of 64 seconds (linear regression) and 18 seconds (gradient boosting).}
        \resumeItem{Analyzed and visualized the dispatch-policy experiments: four Double DQN (RL) policies against the hand-built rule baseline (ETA\_MIN) over 5 seeds and 100k simulated trips; all four beat ETA\_MIN on reward (+5.2\% to +10.7\%), the best per-mine policy ($-267.4$) beat even a greedy oracle ($-272.3$), with 13.5\% lower queue share.}
        \resumeItem{Built a live operations dashboard that turns model outputs into a monitoring view of fleet KPIs and forecasted congestion hotspots on the mine road graph, supporting dispatch decisions.}
      \resumeItemListEnd
  \resumeSubHeadingListEnd

%-----------PROJECTS-----------
\section{Projects}
    \resumeSubHeadingListStart
      \resumeProjectHeading
          {\textbf{Screen Time \& Sleep: Multi-Scale Statistical Study} $|$ \emph{Python (statsmodels), SQL (SQLite)} $|$ \href{https://github.com/VrundaJoshi25/screen-time-sleep-analysis.git}{\faGithub}}{Winter 2026}
          \resumeItemListStart
            \resumeItem{Designed a two-scale study of whether bedtime screen use predicts sleep duration: a self-conducted stratified pilot (n = 30 students, paired exam and regular weeks) plus an independent 996-student public dataset (Mendeley Data, CC BY 4.0).}
            \resumeItem{Analyzed with Spearman and Kendall correlations, HC3-robust OLS, and ordinal-logistic checks, validated by residual, Q-Q, and Cook's Distance diagnostics plus an MCAR/MAR/MNAR sensitivity simulation; built a reproducible pipeline with a SQLite SQL layer (GROUP BY, CTEs, window functions), pytest on GitHub Actions CI, and pinned dependencies.}
            \resumeItem{Both scales rejected the claim that screens cost sleep: pilot exam-week slope p = 0.58 with the 1.15 h exam-week sleep loss pinned on workload, and the 996-student association (+0.38 h per use level) halved to +0.19 h after caffeine, stress, and activity adjustment, consistent with schedule-flexibility confounding via a causal DAG; pre-registered n $\geq$ 100 follow-up (87\% power, phone-log capture) in collection.}
          \resumeItemListEnd
      \resumeProjectHeading
          {\textbf{Heart Disease Prevalence: Interactive Data Visualization} $|$ \emph{scikit-learn, SHAP, Tableau} $|$ \href{https://github.com/VrundaJoshi25/heart-disease-tableau-dashboard.git}{\faGithub}}{Autumn 2026}
          \resumeItemListStart
            \resumeItem{Analyzed a public clinical dataset (Kaggle/UCI); audit found 723 exact duplicate patient records in the raw 1,025-row file and removed them (302 unique patients) to prevent train/test leakage, then trained a Random Forest for feature importance, validated on a stratified 61-patient held-out test set (accuracy 0.787, ROC-AUC 0.865).}
            \resumeItem{Ranked feature importance three ways (impurity-based, permutation on held-out data, SHAP for direction), then separately within male and female subgroups with 500-iteration bootstrap intervals; the ranking shifts descriptively by gender (men: vessel count and ST depression; women: exercise-induced angina and thalassemia test), reported as suggestive because every interval overlaps at n = 302.}
            \resumeItem{Built an interactive Tableau dashboard (gender-split importance, prevalence heatmap, grouped bars; Gender, Age group, and Chest-pain type filters), presented to the class with the course instructor as project jury.}
          \resumeItemListEnd
    \resumeSubHeadingListEnd

%-----------TECHNICAL SKILLS-----------
\section{Technical Skills}
 \begin{itemize}[leftmargin=0in, label={}]
    \small{\item{
     \textbf{Languages}{: Python, SQL} \\
     \textbf{Libraries \& Tools}{: pandas, NumPy, scikit-learn, statsmodels, Matplotlib, Tableau, Excel, Git/GitHub} \\
     \textbf{Concepts}{: Statistical Inference, Hypothesis Testing, Regression Modeling, Data Cleaning \& Visualization}
    }}
 \end{itemize}

%-----------EDUCATION-----------
\section{Education}
  \resumeSubHeadingListStart
    \resumeSubheading
      {Dhirubhai Ambani University}{2025 - Present}
      {M.Sc. Data Science}{Gandhinagar, Gujarat}
    \resumeSubheading
      {Sir P.T. Sarvajanik College of Science}{2022 - 2025}
      {B.Sc. in Mathematics}{Surat, Gujarat}
  \resumeSubHeadingListEnd

%-------------------------------------------
\end{document}
